% persistence-diagram.tex — example persistence diagram (birth vs death)
% Include via % persistence-diagram.tex — example persistence diagram (birth vs death)
% Include via % persistence-diagram.tex — example persistence diagram (birth vs death)
% Include via % persistence-diagram.tex — example persistence diagram (birth vs death)
% Include via \input{figures/persistence-diagram} from overview.tex

\begin{figure}[H]
\centering
\begin{tikzpicture}[
    >=Stealth,
    h0point/.style={circle, fill=blue!70!black, inner sep=2.5pt},
    h1point/.style={diamond, fill=red!70!black, inner sep=2.5pt},
    annotation/.style={font=\small, anchor=west},
]

% Axes
\draw[->, thick] (-0.3, 0) -- (5.5, 0)
    node[anchor=north, font=\small] {birth};
\draw[->, thick] (0, -0.3) -- (0, 5.5)
    node[anchor=east, font=\small] {death};

% Diagonal y = x
\draw[gray, dashed, thin] (0, 0) -- (5.2, 5.2);
\node[gray, font=\footnotesize, rotate=45, anchor=south] at (4.2, 4.2) {$y = x$};

% --- H_0 features (blue circles) ---
% Feature 1: born at 0, dies at 0.9 (short-lived component merge)
\node[h0point] (h0a) at (0, 0.9) {};

% Feature 2: born at 0, dies at 1.2
\node[h0point] (h0b) at (0, 1.2) {};

% Feature 3: born at 0, never dies (essential feature) — plot near top
\node[h0point] (h0c) at (0, 5.0) {};
\draw[blue!70!black, thick, dashed] (h0c) -- (0, 5.4);

% --- H_1 features (red diamonds) ---
% Feature 4: born at 0.9, dies at 1.2 (short-lived loop)
\node[h1point] (h1a) at (0.9, 1.2) {};

% Feature 5: born at 0.7, dies at 2.0 (longest-lived loop — the one we mark)
\node[h1point] (h1b) at (0.7, 2.0) {};

% --- Annotate longest-living feature ---
\draw[red!70!black, thin, ->] (1.8, 2.6) -- ($(h1b) + (0.15, 0.15)$);
\node[annotation, red!70!black] at (1.8, 2.6) {longest-lived $H_1$ feature};

% --- Annotate essential component ---
\draw[blue!70!black, thin, ->] (1.2, 4.8) -- ($(h0c) + (0.15, 0.0)$);
\node[annotation, blue!70!black] at (1.2, 4.8) {essential component ($\infty$)};

% --- Persistence as vertical distance from diagonal ---
% Show for the longest-lived H_1 feature
\draw[red!50!black, thin, dotted] (0.7, 0.7) -- (h1b);
\node[font=\scriptsize, red!50!black, anchor=west] at (0.85, 1.35) {persistence};

% Axis tick marks
\foreach \x in {1, 2, 3, 4, 5} {
    \draw (\x, 0.07) -- (\x, -0.07) node[below, font=\scriptsize] {\x};
    \draw (0.07, \x) -- (-0.07, \x) node[left, font=\scriptsize] {\x};
}

% Legend
\node[h0point, label={[font=\footnotesize]right:$H_0$ (components)}]
    at (3.5, 4.8) {};
\node[h1point, label={[font=\footnotesize]right:$H_1$ (loops)}]
    at (3.5, 4.3) {};

\end{tikzpicture}
\caption{%
    A persistence diagram plots each topological feature as a point at
    coordinates (birth, death). Features far from the diagonal $y = x$
    persist across a wide range of $\varepsilon$ and represent robust
    structure; features near the diagonal are short-lived noise. $H_0$
    features (blue circles) track connected components; $H_1$ features (red
    diamonds) track loops. The dashed extension marks an essential component
    that never dies.%
}
\label{fig:persistence-diagram}
\end{figure}
 from overview.tex

\begin{figure}[H]
\centering
\begin{tikzpicture}[
    >=Stealth,
    h0point/.style={circle, fill=blue!70!black, inner sep=2.5pt},
    h1point/.style={diamond, fill=red!70!black, inner sep=2.5pt},
    annotation/.style={font=\small, anchor=west},
]

% Axes
\draw[->, thick] (-0.3, 0) -- (5.5, 0)
    node[anchor=north, font=\small] {birth};
\draw[->, thick] (0, -0.3) -- (0, 5.5)
    node[anchor=east, font=\small] {death};

% Diagonal y = x
\draw[gray, dashed, thin] (0, 0) -- (5.2, 5.2);
\node[gray, font=\footnotesize, rotate=45, anchor=south] at (4.2, 4.2) {$y = x$};

% --- H_0 features (blue circles) ---
% Feature 1: born at 0, dies at 0.9 (short-lived component merge)
\node[h0point] (h0a) at (0, 0.9) {};

% Feature 2: born at 0, dies at 1.2
\node[h0point] (h0b) at (0, 1.2) {};

% Feature 3: born at 0, never dies (essential feature) — plot near top
\node[h0point] (h0c) at (0, 5.0) {};
\draw[blue!70!black, thick, dashed] (h0c) -- (0, 5.4);

% --- H_1 features (red diamonds) ---
% Feature 4: born at 0.9, dies at 1.2 (short-lived loop)
\node[h1point] (h1a) at (0.9, 1.2) {};

% Feature 5: born at 0.7, dies at 2.0 (longest-lived loop — the one we mark)
\node[h1point] (h1b) at (0.7, 2.0) {};

% --- Annotate longest-living feature ---
\draw[red!70!black, thin, ->] (1.8, 2.6) -- ($(h1b) + (0.15, 0.15)$);
\node[annotation, red!70!black] at (1.8, 2.6) {longest-lived $H_1$ feature};

% --- Annotate essential component ---
\draw[blue!70!black, thin, ->] (1.2, 4.8) -- ($(h0c) + (0.15, 0.0)$);
\node[annotation, blue!70!black] at (1.2, 4.8) {essential component ($\infty$)};

% --- Persistence as vertical distance from diagonal ---
% Show for the longest-lived H_1 feature
\draw[red!50!black, thin, dotted] (0.7, 0.7) -- (h1b);
\node[font=\scriptsize, red!50!black, anchor=west] at (0.85, 1.35) {persistence};

% Axis tick marks
\foreach \x in {1, 2, 3, 4, 5} {
    \draw (\x, 0.07) -- (\x, -0.07) node[below, font=\scriptsize] {\x};
    \draw (0.07, \x) -- (-0.07, \x) node[left, font=\scriptsize] {\x};
}

% Legend
\node[h0point, label={[font=\footnotesize]right:$H_0$ (components)}]
    at (3.5, 4.8) {};
\node[h1point, label={[font=\footnotesize]right:$H_1$ (loops)}]
    at (3.5, 4.3) {};

\end{tikzpicture}
\caption{%
    A persistence diagram plots each topological feature as a point at
    coordinates (birth, death). Features far from the diagonal $y = x$
    persist across a wide range of $\varepsilon$ and represent robust
    structure; features near the diagonal are short-lived noise. $H_0$
    features (blue circles) track connected components; $H_1$ features (red
    diamonds) track loops. The dashed extension marks an essential component
    that never dies.%
}
\label{fig:persistence-diagram}
\end{figure}
 from overview.tex

\begin{figure}[H]
\centering
\begin{tikzpicture}[
    >=Stealth,
    h0point/.style={circle, fill=blue!70!black, inner sep=2.5pt},
    h1point/.style={diamond, fill=red!70!black, inner sep=2.5pt},
    annotation/.style={font=\small, anchor=west},
]

% Axes
\draw[->, thick] (-0.3, 0) -- (5.5, 0)
    node[anchor=north, font=\small] {birth};
\draw[->, thick] (0, -0.3) -- (0, 5.5)
    node[anchor=east, font=\small] {death};

% Diagonal y = x
\draw[gray, dashed, thin] (0, 0) -- (5.2, 5.2);
\node[gray, font=\footnotesize, rotate=45, anchor=south] at (4.2, 4.2) {$y = x$};

% --- H_0 features (blue circles) ---
% Feature 1: born at 0, dies at 0.9 (short-lived component merge)
\node[h0point] (h0a) at (0, 0.9) {};

% Feature 2: born at 0, dies at 1.2
\node[h0point] (h0b) at (0, 1.2) {};

% Feature 3: born at 0, never dies (essential feature) — plot near top
\node[h0point] (h0c) at (0, 5.0) {};
\draw[blue!70!black, thick, dashed] (h0c) -- (0, 5.4);

% --- H_1 features (red diamonds) ---
% Feature 4: born at 0.9, dies at 1.2 (short-lived loop)
\node[h1point] (h1a) at (0.9, 1.2) {};

% Feature 5: born at 0.7, dies at 2.0 (longest-lived loop — the one we mark)
\node[h1point] (h1b) at (0.7, 2.0) {};

% --- Annotate longest-living feature ---
\draw[red!70!black, thin, ->] (1.8, 2.6) -- ($(h1b) + (0.15, 0.15)$);
\node[annotation, red!70!black] at (1.8, 2.6) {longest-lived $H_1$ feature};

% --- Annotate essential component ---
\draw[blue!70!black, thin, ->] (1.2, 4.8) -- ($(h0c) + (0.15, 0.0)$);
\node[annotation, blue!70!black] at (1.2, 4.8) {essential component ($\infty$)};

% --- Persistence as vertical distance from diagonal ---
% Show for the longest-lived H_1 feature
\draw[red!50!black, thin, dotted] (0.7, 0.7) -- (h1b);
\node[font=\scriptsize, red!50!black, anchor=west] at (0.85, 1.35) {persistence};

% Axis tick marks
\foreach \x in {1, 2, 3, 4, 5} {
    \draw (\x, 0.07) -- (\x, -0.07) node[below, font=\scriptsize] {\x};
    \draw (0.07, \x) -- (-0.07, \x) node[left, font=\scriptsize] {\x};
}

% Legend
\node[h0point, label={[font=\footnotesize]right:$H_0$ (components)}]
    at (3.5, 4.8) {};
\node[h1point, label={[font=\footnotesize]right:$H_1$ (loops)}]
    at (3.5, 4.3) {};

\end{tikzpicture}
\caption{%
    A persistence diagram plots each topological feature as a point at
    coordinates (birth, death). Features far from the diagonal $y = x$
    persist across a wide range of $\varepsilon$ and represent robust
    structure; features near the diagonal are short-lived noise. $H_0$
    features (blue circles) track connected components; $H_1$ features (red
    diamonds) track loops. The dashed extension marks an essential component
    that never dies.%
}
\label{fig:persistence-diagram}
\end{figure}
 from overview.tex

\begin{figure}[H]
\centering
\begin{tikzpicture}[
    >=Stealth,
    h0point/.style={circle, fill=blue!70!black, inner sep=2.5pt},
    h1point/.style={diamond, fill=red!70!black, inner sep=2.5pt},
    annotation/.style={font=\small, anchor=west},
]

% Axes
\draw[->, thick] (-0.3, 0) -- (5.5, 0)
    node[anchor=north, font=\small] {birth};
\draw[->, thick] (0, -0.3) -- (0, 5.5)
    node[anchor=east, font=\small] {death};

% Diagonal y = x
\draw[gray, dashed, thin] (0, 0) -- (5.2, 5.2);
\node[gray, font=\footnotesize, rotate=45, anchor=south] at (4.2, 4.2) {$y = x$};

% --- H_0 features (blue circles) ---
% Feature 1: born at 0, dies at 0.9 (short-lived component merge)
\node[h0point] (h0a) at (0, 0.9) {};

% Feature 2: born at 0, dies at 1.2
\node[h0point] (h0b) at (0, 1.2) {};

% Feature 3: born at 0, never dies (essential feature) — plot near top
\node[h0point] (h0c) at (0, 5.0) {};
\draw[blue!70!black, thick, dashed] (h0c) -- (0, 5.4);

% --- H_1 features (red diamonds) ---
% Feature 4: born at 0.9, dies at 1.2 (short-lived loop)
\node[h1point] (h1a) at (0.9, 1.2) {};

% Feature 5: born at 0.7, dies at 2.0 (longest-lived loop — the one we mark)
\node[h1point] (h1b) at (0.7, 2.0) {};

% --- Annotate longest-living feature ---
\draw[red!70!black, thin, ->] (1.8, 2.6) -- ($(h1b) + (0.15, 0.15)$);
\node[annotation, red!70!black] at (1.8, 2.6) {longest-lived $H_1$ feature};

% --- Annotate essential component ---
\draw[blue!70!black, thin, ->] (1.2, 4.8) -- ($(h0c) + (0.15, 0.0)$);
\node[annotation, blue!70!black] at (1.2, 4.8) {essential component ($\infty$)};

% --- Persistence as vertical distance from diagonal ---
% Show for the longest-lived H_1 feature
\draw[red!50!black, thin, dotted] (0.7, 0.7) -- (h1b);
\node[font=\scriptsize, red!50!black, anchor=west] at (0.85, 1.35) {persistence};

% Axis tick marks
\foreach \x in {1, 2, 3, 4, 5} {
    \draw (\x, 0.07) -- (\x, -0.07) node[below, font=\scriptsize] {\x};
    \draw (0.07, \x) -- (-0.07, \x) node[left, font=\scriptsize] {\x};
}

% Legend
\node[h0point, label={[font=\footnotesize]right:$H_0$ (components)}]
    at (3.5, 4.8) {};
\node[h1point, label={[font=\footnotesize]right:$H_1$ (loops)}]
    at (3.5, 4.3) {};

\end{tikzpicture}
\caption{%
    A persistence diagram plots each topological feature as a point at
    coordinates (birth, death). Features far from the diagonal $y = x$
    persist across a wide range of $\varepsilon$ and represent robust
    structure; features near the diagonal are short-lived noise. $H_0$
    features (blue circles) track connected components; $H_1$ features (red
    diamonds) track loops. The dashed extension marks an essential component
    that never dies.%
}
\label{fig:persistence-diagram}
\end{figure}
